\chapter{DIDs: fundamentos, estrutura e resolução}
\label{chap:dids}

\section{Objetivos de aprendizagem}
Ao final deste capítulo, o aluno deve ser capaz de:
\begin{itemize}[leftmargin=*]
  \item explicar o que é um \textit{Decentralized Identifier} (DID) e por que ele difere de identificadores centralizados;
  \item descrever os elementos que compõem um DID Document;
  \item interpretar um exemplo completo de DID Document conforme o padrão;
  \item explicar o processo de resolução de DID e seus principais atores;
  \item distinguir resolução de DID, dereferenciamento e acesso a serviços vinculados.
\end{itemize}
\section{O que é um DID}
Um DID (\textit{Decentralized Identifier}) é um identificador global, controlado criptograficamente, que não depende de uma autoridade central única para sua gestão. De acordo com o padrão DID Core, publicado pelo \textit{World Wide Web Consortium} (W3C), um DID identifica um sujeito --- pessoa, organização, dispositivo, agente de software ou outro recurso --- e aponta para um DID Document com metadados necessários para verificação e interação \citep{w3c_did_core}.

A forma geral de um DID é:
\begin{quote}
\texttt{did:<método>:<identificador-específico>}
\end{quote}

Exemplos:
\begin{itemize}[leftmargin=*]
  \item \texttt{did:web:example.com}
  \item \texttt{did:key:z6Mk...}
  \item \texttt{did:example:123456789abcdefghi}
\end{itemize}

Cada método DID define como o identificador é criado, resolvido, atualizado e, quando aplicável, desativado. Por isso, a semântica operacional de um DID sempre depende do método associado.

\section{Por que usar DIDs}
Identificadores tradicionais costumam depender de diretórios centralizados, provedores de identidade ou bancos de dados administrados por uma única organização. DIDs foram propostos para viabilizar identidades mais portáveis, verificáveis e interoperáveis.

Entre as vantagens arquiteturais mais importantes, destacam-se:
\begin{itemize}[leftmargin=*]
  \item \textbf{Controle criptográfico}: a autenticação e a prova de controle dependem de chaves e assinaturas, e não apenas de cadastro em base central;
  \item \textbf{Portabilidade}: o identificador pode ser reutilizado entre sistemas diferentes, desde que compartilhem o método e as regras de resolução;
  \item \textbf{Verificabilidade}: chaves, serviços e vínculos podem ser auditados com base em documentos e evidências técnicas;
  \item \textbf{Separação entre identidade e aplicação}: o DID representa a identidade; as regras de negócio ficam em camadas superiores.
\end{itemize}

\section{DID Document}
O DID Document é a representação estruturada dos metadados associados a um DID. Em implementações práticas, ele costuma ser serializado em \textit{JavaScript Object Notation} (JSON). Ele informa como um sujeito pode ser autenticado, quais métodos de verificação são válidos e quais serviços podem ser descobertos por meio daquela identidade.

Nem todo DID Document terá exatamente os mesmos campos. O conteúdo depende:
\begin{itemize}[leftmargin=*]
  \item do método DID;
  \item do domínio da aplicação;
  \item dos mecanismos de autenticação e prova adotados;
  \item das extensões e registries utilizados.
\end{itemize}

\section{Exemplo completo de DID Document}
O exemplo a seguir ilustra um DID Document didático contendo os campos mais comuns do padrão:

Antes de examinar o JSON, vale notar três siglas presentes no exemplo: \textit{JSON Web Key} (JWK), usado no campo \texttt{publicKeyJwk}; \textit{JSON Web Signature} (JWS), referenciado no contexto \texttt{jws-2020}; e \textit{Elliptic Curve} (EC), indicado no campo \texttt{kty}.

\begin{verbatim}
{
  "@context": [
    "https://www.w3.org/ns/did/v1",
    "https://w3id.org/security/suites/jws-2020/v1"
  ],
  "id": "did:example:123456789abcdefghi",
  "controller": "did:example:123456789abcdefghi",
  "verificationMethod": [
    {
      "id": "did:example:123456789abcdefghi#key-1",
      "type": "JsonWebKey2020",
      "controller": "did:example:123456789abcdefghi",
      "publicKeyJwk": {
        "kty": "EC",
        "crv": "P-256",
        "x": "f83OJ3D2xF4r5n6Q7s8t9u0v1w2y3z4A5B6C7D8E9F0",
        "y": "x_FEzRu9h2QvW0Ykq1c2d3e4f5g6h7i8j9kLmNoPqRs"
      }
    }
  ],
  "authentication": [
    "did:example:123456789abcdefghi#key-1"
  ],
  "assertionMethod": [
    "did:example:123456789abcdefghi#key-1"
  ],
  "keyAgreement": [
    "did:example:123456789abcdefghi#key-1"
  ],
  "service": [
    {
      "id": "did:example:123456789abcdefghi#messaging",
      "type": "MessagingService",
      "serviceEndpoint": "https://example.org/messages"
    }
  ]
}
\end{verbatim}

\section{Explicação dos campos do exemplo}
\begin{itemize}[leftmargin=*]
  \item \texttt{@context}: define os contextos semânticos usados na interpretação do documento. O primeiro contexto normalmente aponta para o núcleo do padrão DID. No exemplo, o segundo contexto referencia a suíte \textit{JSON Web Signature} (JWS) 2020.
  \item \texttt{id}: é o próprio DID que este documento descreve. Deve coincidir com o identificador resolvido.
  \item \texttt{controller}: informa quem controla o DID ou o recurso descrito. Em muitos casos, o controlador coincide com o próprio \texttt{id}, mas isso não é obrigatório.
  \item \texttt{verificationMethod}: lista os métodos de verificação associados ao DID. Cada entrada descreve uma chave ou mecanismo de verificação.
  \item \texttt{verificationMethod[].id}: identificador único do método de verificação. Frequentemente usa um fragmento, como \texttt{\#key-1}.
  \item \texttt{verificationMethod[].type}: tipo do método de verificação, conforme registries e suites criptográficas aceitas.
  \item \texttt{verificationMethod[].controller}: entidade controladora daquele método específico.
  \item \texttt{publicKeyJwk}: representação da chave pública em formato \textit{JSON Web Key} (JWK). Em outros documentos, a chave pode aparecer em formatos alternativos, como \texttt{publicKeyMultibase}. No exemplo, o valor \texttt{EC} em \texttt{kty} indica o uso de \textit{Elliptic Curve}.
  \item \texttt{authentication}: lista quais métodos podem ser usados para autenticação do sujeito.
  \item \texttt{assertionMethod}: lista os métodos autorizados para emitir ou assinar afirmações, como credenciais e declarações verificáveis.
  \item \texttt{keyAgreement}: lista métodos aptos a estabelecer segredos compartilhados, em cenários de troca segura de mensagens.
  \item \texttt{service}: lista pontos de serviço vinculados ao DID, como endpoints de mensageria, descoberta, resolução auxiliar ou interfaces de programação de aplicações (APIs) específicas.
  \item \texttt{serviceEndpoint}: endereço do serviço associado. Não é, por si só, prova de autenticidade; é apenas um vínculo declarado pelo documento.
\end{itemize}

\section{Campos mínimos e extensões}
Na prática, um DID Document pode ser mais enxuto ou mais rico do que o exemplo anterior. Há documentos que incluem apenas:
\begin{itemize}[leftmargin=*]
  \item \texttt{@context};
  \item \texttt{id};
  \item um ou mais \texttt{verificationMethod};
  \item uma forma de autenticação.
\end{itemize}

Por outro lado, aplicações específicas podem acrescentar:
\begin{itemize}[leftmargin=*]
  \item serviços personalizados;
  \item chaves segregadas por propósito;
  \item metadados operacionais;
  \item vínculos com infraestrutura de armazenamento distribuído ou registradores externos.
\end{itemize}

\section{Resolução de DID}
A resolução de DID é o processo de obter, a partir de um DID, o DID Document correspondente e os metadados associados a essa operação. Em termos conceituais, resolver um DID significa responder à pergunta: \textit{``qual é o documento atual e verificável que descreve esta identidade?''}

Esse processo envolve, em geral, quatro atores:
\begin{itemize}[leftmargin=*]
  \item \textbf{Cliente}: aplicação que precisa do DID Document;
  \item \textbf{Resolver}: componente que interpreta o método DID e coordena a busca;
  \item \textbf{Driver ou adaptador do método}: implementação específica de cada método;
  \item \textbf{Fonte de verdade}: infraestrutura onde o estado do DID está registrado, como \textit{Hypertext Transfer Protocol} (HTTP), blockchain, \textit{Distributed Hash Table} (DHT), banco distribuído ou outro registro verificável.
\end{itemize}

\section{Saídas da resolução}
Em uma implementação completa, a resolução não devolve apenas o DID Document. Ela costuma produzir três grupos de informação:
\begin{itemize}[leftmargin=*]
  \item \textbf{DID Document}: o conteúdo resolvido;
  \item \textbf{Metadados do documento}: informações como versão, atualização, desativação ou equivalências;
  \item \textbf{Metadados da resolução}: informações sobre o próprio processo, incluindo erros, método usado, cache e opções de resolução.
\end{itemize}

\section{Fluxo genérico de resolução}
Embora cada método DID tenha seu próprio comportamento, o fluxo genérico costuma seguir a sequência abaixo:
\begin{enumerate}[leftmargin=*]
  \item o cliente fornece um DID ao resolvedor;
  \item o resolvedor identifica o método, por exemplo \texttt{web}, \texttt{key} ou outro;
  \item o adaptador daquele método aplica as regras de resolução correspondentes;
  \item o sistema consulta a fonte de verdade do método;
  \item o documento obtido é validado e normalizado;
  \item o resultado é devolvido ao cliente com documento e metadados.
\end{enumerate}

\section{Exemplos de resolução por método}
\begin{itemize}[leftmargin=*]
  \item \textbf{\texttt{did:web}}: a resolução é feita via \textit{Hypertext Transfer Protocol Secure} (HTTPS), normalmente recuperando um arquivo como \texttt{/.well-known/did.json} ou caminho equivalente no domínio.
  \item \textbf{\texttt{did:key}}: a resolução é local e determinística, pois a chave pública já está embutida no próprio identificador.
  \item \textbf{Métodos ancorados em blockchain}: a resolução exige consulta ao registro on-chain, leitura de eventos, estados ou apontadores para dados externos.
\end{itemize}

\section{Resolução em arquiteturas com registro verificável e IPFS}
Em arquiteturas didáticas ou experimentais, é comum separar:
\begin{itemize}[leftmargin=*]
  \item um \textbf{registro verificável}, responsável por manter o vínculo entre DID e uma referência estável;
  \item um \textbf{armazenamento distribuído}, como o \textit{InterPlanetary File System} (IPFS), responsável por hospedar o documento completo.
\end{itemize}

Um fluxo típico é:
\begin{enumerate}[leftmargin=*]
  \item resolver o DID no registro verificável;
  \item obter um ponteiro, como um hash, um \textit{Content Identifier} (CID) ou um identificador de versão;
  \item recuperar o documento na camada de armazenamento;
  \item validar se o documento recuperado é coerente com o vínculo registrado.
\end{enumerate}

Essa abordagem ajuda a separar governança de identidade, auditabilidade e armazenamento por conteúdo.

\section{Resolução versus dereferenciamento}
É importante não confundir resolução com dereferenciamento.
\begin{itemize}[leftmargin=*]
  \item \textbf{Resolução}: retorna o DID Document associado a um DID.
  \item \textbf{Dereferenciamento}: resolve um recurso específico dentro desse universo, como um fragmento \texttt{\#key-1} ou um serviço particular.
\end{itemize}

Exemplo:
\begin{itemize}[leftmargin=*]
  \item resolver \texttt{did:example:123456789abcdefghi} significa obter o DID Document completo;
  \item dereferenciar \texttt{did:example:123456789abcdefghi\#key-1} significa obter o método de verificação específico apontado pelo fragmento.
\end{itemize}

\section{Cuidados técnicos no processo de resolução}
Ao projetar ou implementar resolução de DID, alguns cuidados são essenciais:
\begin{itemize}[leftmargin=*]
  \item validar a autenticidade da fonte de verdade;
  \item tratar indisponibilidade, timeout e inconsistência entre registro e documento;
  \item registrar metadados de erro e rastreabilidade da resolução;
  \item considerar cache com critério, para não servir documentos desatualizados;
  \item proteger informações sensíveis que não devam ser publicadas em documentos resolvíveis.
\end{itemize}

\section{Boas práticas de modelagem}
\begin{itemize}[leftmargin=*]
  \item manter o DID Document focado em dados necessários para verificação e descoberta;
  \item separar dados sensíveis de metadados públicos;
  \item escolher formatos de chave e serviços alinhados ao ecossistema do método DID;
  \item documentar critérios de atualização, revogação e desativação.
\end{itemize}

\section{Exercícios sugeridos}
\begin{enumerate}[leftmargin=*]
  \item comparar os métodos \texttt{did:key}, \texttt{did:web} e um método ancorado em blockchain;
  \item identificar, no exemplo deste capítulo, quais campos são voltados para autenticação, quais são voltados para assinatura de afirmações e quais são voltados para descoberta de serviços;
  \item desenhar um fluxo de resolução para uma arquitetura que use contrato inteligente e IPFS;
  \item explicar, com um diagrama simples, a diferença entre resolução e dereferenciamento.
\end{enumerate}

\section{Leituras recomendadas}
\begin{itemize}[leftmargin=*]
  \item W3C DID Core, norma principal sobre modelo de dados e semântica \citep{w3c_did_core};
  \item W3C DID Specification Registries, para métodos, tipos e registries de apoio \citep{w3c_did_spec_registries};
  \item DID Resolution, para aprofundamento do processo de resolução e metadados \citep{w3c_did_resolution};
  \item \textit{Request for Comments} (RFC) 3986, para sintaxe geral de identificadores \textit{Uniform Resource Identifier} (URI) \citep{rfc3986};
  \item W3C Verifiable Credentials Data Model, para integração entre identidade e credenciais verificáveis \citep{w3c_vc_data_model}.
\end{itemize}

\section{Rodando na prática}
O objetivo desta prática é executar um resolvedor mínimo de \texttt{did:web}, publicar localmente um \textit{DID Document} e observar, passo a passo, como um DID é convertido em um recurso resolvível.

\subsection{Pré-requisitos}
\begin{itemize}[leftmargin=*]
  \item Node.js 18 ou superior;
  \item terminal com suporte aos comandos do Node;
  \item diretório \texttt{docs/book/examples/cap1-did-resolution} disponível localmente.
\end{itemize}

\subsection{Arquivos do laboratório}
O mini projeto deste capítulo contém:
\begin{verbatim}
did.json
server.mjs
resolve.mjs
\end{verbatim}

\subsection{Executando}
No diretório do laboratório, abra dois terminais. No primeiro:
\begin{verbatim}
node server.mjs
\end{verbatim}

No segundo:
\begin{verbatim}
node resolve.mjs did:web:localhost%3A8081:users:alice
\end{verbatim}

\subsection{O que observar}
\begin{itemize}[leftmargin=*]
  \item a regra de transformação do identificador \texttt{did:web} em URL;
  \item a recuperação do documento \texttt{did.json} por HTTP;
  \item a estrutura retornada, com \texttt{id}, \texttt{verificationMethod}, \texttt{authentication} e \texttt{service};
  \item a diferença entre o identificador abstrato e o recurso concreto acessado durante a resolução.
\end{itemize}

\subsection{Perguntas de análise}
\begin{itemize}[leftmargin=*]
  \item Por que o laboratório usa HTTP em vez de HTTPS?
  \item Qual é a diferença entre resolver um DID e dereferenciar um fragmento como \texttt{\#key-1}?
  \item O que aconteceria se o arquivo \texttt{did.json} não contivesse o campo \texttt{id}?
  \item Como esse mecanismo de resolução se compara a uma consulta em diretório centralizado?
\end{itemize}

\subsection{Extensão sugerida}
Altere o arquivo \texttt{did.json} para:
\begin{itemize}[leftmargin=*]
  \item adicionar um novo \texttt{serviceEndpoint};
  \item trocar o identificador do controlador;
  \item incluir uma segunda chave em \texttt{verificationMethod}.
\end{itemize}

Depois repita a execução e verifique como a mudança afeta a saída do resolvedor.

\subsection{Localização do exemplo}
Todos os arquivos usados nesta prática estão em:
\begin{verbatim}
docs/book/examples/cap1-did-resolution
\end{verbatim}

\vspace{1em}
\noindent\rule{\textwidth}{0.4pt}
\vspace{0.5em}
\noindent\textbf{Transição para o próximo capítulo.}
Até aqui, o foco esteve na camada lógica da identidade: como um DID é estruturado, publicado e resolvido. No entanto, em um sistema de controle de acesso físico como o DIDGuard, essa identidade precisa estar associada a uma credencial tangível --- um cartão, tag ou dispositivo que o portador apresenta a um leitor. O próximo capítulo explora essa camada física, investigando como as tags RFID MIFARE Classic organizam sua memória, autenticam setores e --- sobretudo --- por que suas proteções nativas não são suficientes quando usadas isoladamente.

